% \iffalse meta-comment
%
% Transformed from classpack.xml by ClassPack db2dtx.xsl
% version 1.19 (2020-05-19) on Wednesday 20 May 2020 at 14:32:36
%
% classpack.dtx is copyright © 2012-2020 by Peter Flynn <peter@silmaril.ie>
%
% This work may be distributed and/or modified under the
% conditions of the LaTeX Project Public License, either
% version 1.3 of this license or (at your option) any later
% version. The latest version of this license is in:
%
%     http://www.latex-project.org/lppl.txt
%
% and version 1.3 or later is part of all distributions of
% LaTeX version 2005/12/01 or later.
%
% This work has the LPPL maintenance status ‘maintained’.
% 
% The current maintainer of this work is Peter Flynn <peter@silmaril.ie>
%
% This work consists of the files classpack.dtx and classpack.ins,
% the derived file , 
% and any other ancillary files listed in the MANIFEST.
%
% \fi
% \iffalse
%<package>\NeedsTeXFormat{LaTeX2e}[2015/01/01]
%<package>\ProvidesPackage{classpack}[2020/05/19 v1.19
%<package>  Macros for ClassPack documentation]
%<*driver>
\RequirePackage{fix-cm}% included by default.
\PassOptionsToPackage{svgnames}{xcolor}% xcolor/dox/hyperref implied
\documentclass[12pt]{ltxdoc}
%%
%% Packages added for documentation
%%
\usepackage{dox}% used by default. (0)%
  \makeatletter
  \doxitem[idxtype=attribute]{Attribute}{CPK@attribute}{attributes}
  \makeatother
  \makeatletter
  \doxitem[idxtype=attributevalue]{AttributeValue}{CPK@attributevalue}{attribute values}
  \makeatother
  \makeatletter
  \doxitem[idxtype=class]{Class}{CPK@class}{classes}
  \makeatother
  \makeatletter
  \doxitem[idxtype=colour]{Colour}{CPK@colour}{colours}
  \makeatother
  \makeatletter
  \doxitem[idxtype=counter]{Counter}{CPK@counter}{counters}
  \makeatother
  \makeatletter
  \doxitem[idxtype=DTD]{DTD}{CPK@dtd}{DTDs/Schemas}
  \makeatother
  \makeatletter
  \doxitem[idxtype=element]{Element}{CPK@element}{element types}
  \makeatother
  \makeatletter
  \doxitem[idxtype=entity]{Entity}{CPK@entity}{entities}
  \makeatother
  \makeatletter
  \doxitem[idxtype=error]{Error}{CPK@error}{errors}
  \makeatother
  \makeatletter
  \doxitem[idxtype=field]{Field}{CPK@field}{fields}
  \makeatother
  \makeatletter
  \doxitem[idxtype=file]{File}{CPK@file}{files}
  \makeatother
  \makeatletter
  \doxitem[idxtype=font]{Font}{CPK@font}{fonts}
  \makeatother
  \makeatletter
  \doxitem[idxtype=function]{Function}{CPK@function}{functions}
  \makeatother
  \makeatletter
  \doxitem[idxtype=language]{Language}{CPK@language}{languages}
  \makeatother
  \makeatletter
  \doxitem[macrolike,idxtype=length]{Length}{CPK@length}{lengths}
  \makeatother
  \makeatletter
  \doxitem[idxtype=mode]{Mode}{CPK@mode}{modes}
  \makeatother
  \makeatletter
  \doxitem[idxtype=option]{Option}{CPK@option}{options}
  \makeatother
  \makeatletter
  \doxitem[idxtype=package]{Package}{CPK@package}{packages}
  \makeatother
  \makeatletter
  \doxitem[idxtype=variable]{Variable}{CPK@variable}{variables}
  \makeatother
  \makeatletter
  \doxitem[idxtype=parameter]{Parameter}{CPK@parameter}{parameters}
  \makeatother
  \makeatletter
  \doxitem[macrolike,idxtype=switch]{Switch}{CPK@switch}{switches}
  \makeatother
  \makeatletter
  \doxitem[idxtype=template]{Template}{CPK@template}{templates}
  \makeatother
  \makeatletter
  \doxitem[idxtype=typeface]{Typeface}{CPK@typeface}{typefaces}
  \makeatother
  \makeatletter
  \doxitem[macrolike,idxtype=box]{Box}{CPK@box}{boxes}
  \makeatother
  \newcommand{\LabelFont}[2][\relax]{\strut
    {\fontencoding\encodingdefault
	    \fontfamily{lmtt}\fontseries{lc}#1\selectfont#2}\space}
  \makeatletter
  \let\CPK@macro\macro\let\CPK@endmacro\endmacro
  \makeatother
  \makeatletter
  \let\CPK@environment\environment\let\CPK@endenvironment\endenvironment
  \makeatother
  \makeatletter
  \def\PrintAttributeName#1{\LabelFont{@#1}}
  \makeatother
  \def\PrintAttributeValueName#1{\LabelFont{"#1"}}
  \def\PrintClassName#1{\LabelFont[\fontfamily{lmss}]{#1}}
  \def\PrintColourName#1{\LabelFont[\color{#1}]{#1}}
  \def\PrintCounterName#1{\LabelFont{#1}}
  \def\PrintDTDName#1{\LabelFont{#1}}
  \def\PrintElementName#1{\LabelFont{<#1>}}
  \def\PrintEntityName#1{\LabelFont{\&#1;}}
  \def\PrintEnvironmentName#1{\LabelFont[\fontfamily{lmss}]{#1}}
  \def\PrintErrorName#1{\LabelFont[\color{Red}!]{#1}}
  \def\PrintFunctionName#1{\LabelFont[\bfseries\itshape]{#1}}
  \def\PrintLanguageName#1{\LabelFont{#1}}
  \def\PrintLengthName#1{\LabelFont{#1}}
  \def\PrintMacroName#1{\LabelFont{#1}}
  \def\PrintModeName#1{\LabelFont[\sffamily]{\textlangle#1\textrangle}}
  \def\PrintOptionName#1{\LabelFont[\bfseries]{#1}}
  \def\PrintPackageName#1{\LabelFont[\fontfamily{lmss}]{#1}}
  \def\PrintSwitchName#1{\LabelFont{#1}}
  \def\PrintTemplateName#1{\LabelFont[\bfseries]{#1}}
  \def\PrintVariableName#1{\LabelFont[\ttfamily]{#1}}
  \def\PrintParameterName#1{\LabelFont[\ttfamily]{#1}}
  \def\PrintFieldName#1{\LabelFont[\ttfamily]{#1}}
%% fontenc omit: conflict: fontspec (3)
%% inputenc omit: conflict: fontspec (4)
\usepackage{fontspec}% part/@conformance=xelatex detected. (7)%
\usepackage{mflogo}% used by default. (35)%
\usepackage[greek,british]{babel}% requested by author (40)%
\usepackage[backend=biber,doi=true,
  isbn=true,url=true,uniquename=false,style=apa]{biblatex}% bibliography/@arch=biblatex detected. (41)%
  \makeatletter
  \AtBeginDocument{%
	    \setlength{\bibitemsep}{1ex}%
	    \setlength{\bibnamesep}{1.5\itemsep}%
	    \defbibheading{shortbib}[References]%
	      {\section{#1}}}
	  \@ifpackagewith{babel}{british}{%
	    \DeclareLanguageMapping{british}%
	      {british-apa}}{\relax}
  \makeatother
  \makeatletter
  \providetoggle{blx@skipbiblist}
  \makeatother
\usepackage{csquotes}% bibliography/@arch=biblatex detected. (43)%
\usepackage{array}% use of 'tgroup' detected (46)%
\usepackage{calc}% used by default. (49)%
  \makeatletter
  {\scriptsize
	    \global\advance\@totalleftmargin by1em
	    \global\advance\MacroIndent by.5em}
  \makeatother
\usepackage{ccaption}% used by default. (52)%
  \captionnamefont{\bfseries}
  \captionstyle{\raggedright}
\usepackage[inline]{enumitem}% use of 'variablelist' detected (56)%
  \setlist[description]{style=unboxed}
  \setlist[itemize]{leftmargin=2em}
  \setlist[enumerate]{leftmargin=2em}
  \newlist{inlineenum}{enumerate*}{1}
  \setlist[inlineenum,1]{label=\emph{\alph*}),
	    itemjoin={{; }},itemjoin*={{; and }}}
\usepackage{fancybox}% use of 'note' detected (57)%
\usepackage{fancyvrb}% use of 'bibliography' detected (60)%
\usepackage{relsize}% use of 'acronym' detected (64)%
\usepackage{textcase}% used by default. (67)%
\usepackage{float}% used by default. (69)%
  \renewcommand{\topfraction}{.85}
  \renewcommand{\bottomfraction}{.7}
  \renewcommand{\textfraction}{.15}
  \renewcommand{\floatpagefraction}{.66}
  \renewcommand{\dbltopfraction}{.66}
  \renewcommand{\dblfloatpagefraction}{.66}
  \setcounter{topnumber}{9}
  \setcounter{bottomnumber}{9}
  \setcounter{totalnumber}{20}
  \setcounter{dbltopnumber}{9}
\usepackage[a4paper,left=30mm,top=25mm,
  textwidth=150mm,textheight=225mm,headheight=15pt]{geometry}% used by default. (71)%
\usepackage{graphicx}% part/@conformance=xelatex detected. (76)%
\usepackage{listings}% use of 'programlisting' detected (83)%
  \lstdefinelanguage{dummy}
	    {morekeywords={dummy}}
  \lstdefinelanguage{Makefile}
	  {otherkeywords={.PHONY,.DEFAULT},%
	    morekeywords={PHONY,DEFAULT,shell,ifeq,else,endif},%
	    keywordsprefix={.},%
	    moredelim=[l][\color{Green}]{:},%
	    morecomment=[l]{\#},%
	    moredelim=[s][\color{Blue}]{\$(}{)}%
	  }
  \lstdefinelanguage{DocBook}[]{XML}
	    {morekeywords={abstract,address,affiliation,annotation,arg,
	    author,book,chapter,classname,cmdsynopsis,command,
	    constraintdef,contrib,copyright,cover,date,email,emphasis,
	    envar,filename,firstname,footnote,guibutton,guilabel,
	    guimenu,guimenuitem,guisubmenu,holder,info,itemizedlist,
	    listitem,literal,member,option,orderedlist,orgdiv,orgname,
	    package,para,parameter,part,personname,phrase,procedure,
	    productname,programlisting,quote,refsection,remark,
	    constructorsynopsis,methodparan,modifier,funcparams,olink,
	    bibliography,biblioentry,biblioset,subtitle,artpagenums,
	    volumenum,issuenum,DOCTYPE,SYSTEM,xml:id,releaseinfo,
	    replaceable,revdescription,revhistory,revision,sect1,sect2,
	    sect3,sect4,seg,seglistitem,segmentedlist,segtitle,
	    simplelist,step,surname,systemitem,tag,term,title,uri,
	    userinput,variablelist,varlistentry,wordasword,xref,year,
	    xlink:href}}
	  
  \makeatletter
  \lstdefinelanguage{bash}
	    {morestring=[s]{[]},morekeywords={exit,logout,yes,no,@,
	    password,ssh,URL,cd,dvips,latex,ls,makeindex,man,mkdir,
	    pdflatex,sudo,texconfig,texdoc,updmap,xelatex,biber,
	    latexmk,bibtex}} 
	  
  \makeatother
  \lstdefinelanguage{APA}[]{XML}
	    {morekeywords={TTL}}
	  
  \lstdefinelanguage{OOXML}[]{XML}
	    {morekeywords={w:p,w:pPr,w:pStyle,w:rPr,w:rFonts,
	    w:r,w:t,w:lang}}
	  
  \lstdefinelanguage{SGML}[]{XML}
	    {morekeywords={sec,ttl}}
	  
  \lstdefinelanguage{DTD}[]{XML}
	    {morekeywords={ELEMENT,ENTITY,ATTLIST,CDATA,ID,REQUIRED,
	    IMPLIED,PCDATA}}
	  
  \lstdefinelanguage{Runoff}
	    {morekeywords={h1}}
	  
  \lstdefinelanguage{GML}
	    {morekeywords={h1}}
	  
  \lstdefinelanguage{Scribe}
	    {morekeywords={Heading},morestring=[s]{[]}}
	  
  \lstdefinelanguage{RTF}[]{TeX}
	    {moretexcs={rtf,ansi,deff,adeflang,fonttbl,f,froman,fprq,
	    fcharset,f1,fswiss,falt,fnil,colortbl,red,green,blue,
	    stylesheet,s,snext,nowidctlpar,hyphen,hyphlead,hyphtrail,
	    hyphmax,cf,kerning,dbch,af,langfe,afs,alang,loch,fs,
	    pgndec,pard,plain,qc,sb,sa,keepn,b,ab,rtlch,ltrch,par}}

  \lstdefinelanguage{TEI}[]{XML}
	    {morekeywords={TEI,TEI.2,teiHeader,fileDesc,sourceDesc,
	    titleStmt,title,author,editor,respStmt,resp,name,
	    editionStmt,edition,text,body,publicationStmt,publisher,
	    div,div1,placeName,lg,l,s,cl,phr,w,list,distinct,p,pb,
	    mls,div2,head,num,val,app,lem,rdg,q,sup,uncl,note,
	    DOCTYPE,SYSTEM,xml:id}}[keywords,comments,strings]
	  
  \lstdefinelanguage{XSLT2}[]{XML}
	    {morekeywords={xsl:stylesheet,xsl:transform,
	    xsl:apply-imports,xsl:attribute-set,xsl:decimal-format,
	    xsl:import,xsl:include,xsl:key,xsl:namespace-alias,
	    xsl:output,xsl:param,
	    xsl:preserve-space,xsl:strip-space,xsl:template,
	    xsl:variable,xsl:character-map,xsl:function,
	    xsl:import-schema,xsl:param,xsl:variable,
	    xsl:apply-imports,xsl:apply-templates,xsl:attribute,
	    xsl:call-template,xsl:choose,xsl:comment,xsl:copy,
	    xsl:copy-of,xsl:element,xsl:fallback,xsl:for-each,
	    xsl:if,xsl:message,xsl:number,xsl:otherwise,
	    xsl:processing-instruction,xsl:text,xsl:value-of,
	    xsl:variable,xsl:when,xsl:with-param,xsl:sort,
	    xsl:for-each-group,xsl:next-match,xsl:analyze-string,
	    xsl:namespace,xsl:result-document,xsl:copy,
	    xsl:fallback,xsl:document,xsl:sequence,
	    xsl:matching-substring,xsl:non-matching-substring,
	    xsl:perform-sort,xsl:output-character},
	    alsodigit={-}}
	  
  \lstdefinelanguage{LaTeXe}[LaTeX]{TeX}
	    {morekeywords = {selectlanguage,foreignlanguage,
	    textbrokenbar,textlangle,textrangle,subsection,url,
	    chapter,tableofcontents,part,subsubsection,paragraph,
	    subparagraph,maketitle,setlength,listoffigures,
	    listoftables,color,arraybackslash,includegraphics,
	    textcite,parencite,graphicspath,lstinline,
	    DeclareLanguageMapping,textcolor,definecolor,colorbox,
	    fcolorbox,RequirePackage,PassOptionsToPackage}}
	  
  \lstdefinelanguage{BIBTeX}{
	    morekeywords = {title,author,edition,publisher,year,
	    address},
	    morestring=[b]",
	    }
	  
  \lstdefinelanguage{Email}{
	    morekeywords={From,Subject,To,Date},
	    }
	  
  \lstset{defaultdialect=LaTeXe,frame=single,
	    framesep=.5em,backgroundcolor=\color{AliceBlue},
	    rulecolor=\color{LightSteelBlue},framerule=1pt}
	  
  \lstloadlanguages{LaTeXe,DocBook,XML,XSLT2,bash}
  \lstdefinelanguage{XMLFRAG}{tag=**[s]<>}[html]
  \lstnewenvironment{listingsdoc}
	    {\lstset{language={[LaTeX]TeX}}}
	    {}
  \newcommand\basicdefault[1]{\footnotesize
	    \color{Black}\ttfamily#1}
	  
  \lstset{basicstyle=\basicdefault{\spaceskip.5em}}
  \lstset{literate=
	    {§}{{\S}}1
	    {©}{{\raisebox{.125ex}{\copyright}\enspace}}1
	    {«}{{\guillemotleft}}1
	    {»}{{\guillemotright}}1
	    {Á}{{\'A}}1
	    {Ä}{{\"A}}1
	    {É}{{\'E}}1
	    {Í}{{\'I}}1
	    {Ó}{{\'O}}1
	    {Ö}{{\"O}}1
	    {Ú}{{\'U}}1
	    {Ü}{{\"U}}1
	    {ß}{{\ss}}2
	    {à}{{\`a}}1
	    {á}{{\'a}}1
	    {ä}{{\"a}}1
	    {é}{{\'e}}1
	    {í}{{\'i}}1
	    {ó}{{\'o}}1
	    {ö}{{\"o}}1
	    {ú}{{\'u}}1
	    {ü}{{\"u}}1
	    {¹}{{\textsuperscript1}}1
            {²}{{\textsuperscript2}}1
            {³}{{\textsuperscript3}}1
	    {ı}{{\i}}1
	    {—}{{---}}1
	    {’}{{'}}1
	    {…}{{\dots}}1
	    {➝}{{$leftarrow$}}1
            {⮠}{{$\hookleftarrow$}}1
	    {␣}{{\textvisiblespace}}1,
	    keywordstyle=\color{DarkGreen}\bfseries,
	    identifierstyle=\color{DarkRed},
	    commentstyle=\color{Gray}\upshape,
	    stringstyle=\color{DarkBlue}\upshape,
	    emphstyle=\color{Chocolate}\upshape,
	    showstringspaces=false,
	    columns=fullflexible,
	    keepspaces=true}
\usepackage{makeidx}% used by default. (85)%
  \makeindex
\usepackage{nicefrac}% used by default. (91)%
  \def\textonehalf{\ensuremath{\nicefrac12}}
\usepackage{parskip}% used by default. (93)%
\usepackage{sectsty}% used by default. (96)%
  \allsectionsfont{\sffamily\raggedright}
  \renewcommand*{\descriptionlabel}[1]{\hspace\labelsep
	    \sffamily\bfseries #1}
\usepackage[normalem]{ulem}% use of 'link' detected (103)%
\usepackage{url}% use of 'ulink' detected (104)%
  \AtBeginDocument{\urlstyle{tt}}
\usepackage{varioref}% use of 'xref' detected (107)%
  \vrefwarning
  \labelformat{appendix}{Appendix~#1}
  \makeatletter
  \labelformat{chapter}{\@chapapp~#1}
  \makeatother
  \labelformat{section}{section~#1}
  \labelformat{subsection}{section~#1}
  \labelformat{subsubsection}{section~#1}
  \labelformat{paragraph}{section~#1}
  \labelformat{figure}{Figure~#1}
  \labelformat{table}{Table~#1}
  \labelformat{item}{item~#1}
  \renewcommand{\reftextcurrent}{elsewhere on this
	    page}
  \def\reftextafter{on the
	    \reftextvario{next}{following} page}
\usepackage{xcolor}% used by default. (115)%
  \makeatletter
  \@ifundefined{T}{%
	    \newcommand{\T}[2]{{\fontencoding{T1}%
	      \selectfont#2}}}{}
  \makeatother
\usepackage{fontawesome}% guiicon/@remap detected. (117)%
\usepackage{menukeys}% use of 'guimenu' detected (118)%
  \renewmenumacro{\menu}[>]{roundedmenus}
  \renewmenumacro{\directory}[/]{hyphenatepathswithfolder}
  \renewmenumacro{\keys}{shadowedroundedkeys}
\usepackage{classpack}[2020/05/19]% used by default. (120)%
\usepackage[ragged]{footmisc}% requested by author (9999)%
\usepackage{classpack}[2020/05/19]% added by spec
\addbibresource{classpack.bib}
%
%%
%% Settings for docstrip and ltxdoc 
%%
\EnableCrossrefs
\CodelineIndex
\RecordChanges
\begin{document}\raggedright
  \DocInput{classpack.dtx}
\end{document}
%</driver>
% \fi
%
% \CheckSum{246}
%
% \CharacterTable
%  {Upper-case    \A\B\C\D\E\F\G\H\I\J\K\L\M\N\O\P\Q\R\S\T\U\V\W\X\Y\Z
%   Lower-case    \a\b\c\d\e\f\g\h\i\j\k\l\m\n\o\p\q\r\s\t\u\v\w\x\y\z
%   Digits        \0\1\2\3\4\5\6\7\8\9
%   Exclamation   \!  Double quote  \"  Hash (number) \#
%   Dollar        \$  Percent       \%  Ampersand     \&
%   Acute accent  \'  Left paren    \(  Right paren   \)
%   Asterisk      \*  Plus          \+  Comma         \,
%   Minus         \-  Point         \.  Solidus       \/
%   Colon         \:  Semicolon     \;  Less than     \<
%   Equals        \=  Greater than  \>  Question mark \?
%   Commercial at \@  Left bracket  \[  Backslash     \\
%   Right bracket \]  Circumflex    \^  Underscore    \_
%   Grave accent  \`  Left brace    \{  Vertical bar  \|
%   Right brace   \}  Tilde         \~}
% 
% \changes{v1.19}{2020/05/19}{: 1) Prefixed local variables with CPK; 2) CTAN candidate.}
% \changes{v1.18}{2020/04/01}{: 1) Split project into classpack and classpack-dev; 2) Replaced conformance with YYYY-MM-DD on all date elements in the revision history. This meant a complete regression run after editing all source masters to use this change; 3) Fixed bug in makechapapp where it failed to add the filetype before substringing the full name of the document.}
% \changes{v1.17}{2020/03/19}{: 1) Documentation nearly finished; 2) Added remaining XSLT files and scripts; 3) Recoded the distinction between appendixes in the code part and chapters in the files part..}
% \changes{v1.16}{2019/10/20}{: 1) Documentation ongoing; 2) Added detection of options on url and xcolor; 3) Changed order of packages in prepost.xml to prevent option clashes.}
% \changes{v1.15}{2019/10/10}{: 1) Started massive documentation effort on XSLT; 2) Order of templates changed in XSLT.}
% \changes{v1.14}{2019/07/25}{: 1) Updated the Markdown driver for the README to produce code that displays correctly in the Markdown Viewer 3.6 extension for Chrome and conforms to the CTAN rules.; 2) Added details of ClassPack to the MD file header; 3) Some reorganisation of topics; 4) More annotation-level documentation on the XSLT3 code.}
% \changes{v1.13}{2019/02/18}{: 1) Created commentchar and termcommentchar functions; 2) Untested invocation in sect1 template.}
% \changes{v1.12}{2019/02/17}{: 1) Started tidying up sentinel and fence; 2) Involves rationalising the use of chapter and appendix within part; 3) parameterised a makesentinel, maketermsentinel, and makefence in sect1 but they don't get called because of [2] above; 4) Finished writing the chunk program.}
% \changes{v1.11}{2019/02/13}{: 1) Fixed numerous bugs in handling of listings; 2) Moved DTD annotations to doc as chapter while logic for appendixes not working properly; 3) Finished PE chunking of db2dtx.xsl and doctexbook-master.dtd; 4) Finished writing the chunk program.}
% \changes{v1.10}{2019/01/22}{: 1) Added attributes to programlisting for fragmentary listings; 2) Changed emphasis with role=bold to be plain bold, not bold italic; 3) Changed blockquotes to blue in the PDF, added recognition of xlink:href to give source ack as URI.}
% \changes{v1.09}{2018/07/07}{Skipped the application of headers to appendix files: 1) Files written from the appendix element no longer get the LaTeX headers prepended; 2) The sentinel value for comments in scripts now reflects the double hash; 3) Added the remark element to db2md.xsl (with plain para); 4) Ongoing updates and explanation.}
% \changes{v1.08}{2018/02/07}{Moved RFC2119 warning, and did some minor rewording: 1) Removed RFC2119 warning text from db2dtx.xsl to rfc2119.xml, creating a section to hold it and the bibliography. Rewrote the templates for bibliolist, biblioentry, bibliography, and section/title to accomodate this.; 2) Ongoing updates and explanation.}
% \changes{v1.07}{2018/01/12}{Minor adaptations to quoted chunks of code: Updated to use \XeLaTeX{}.}
% \changes{v1.06}{2017/06/16}{Maintenance release: Rearranged output so that change log and index always get printed.}
% \changes{v1.05}{2017/04/03}{Maintenance: 1) Updated documentation; 2) Tested additional outputs.}
% \changes{v1.04}{2016/12/20}{Reusable code blocks: 1) Added an attribute reuse to the annotation element for use in Appendixes which generate additional \LaTeX{} class or package files, which points at an xml:id attribute on an existing annotation element already used elsewhere so that the same documentation text can be included. Also added an omit attribute to the seglistitem elements in the constraintdef element for class packages which specifies that the relevant package is not required for the additional class or package being generated.; 2) Changed the documentation font to Noto..}
% \changes{v1.03}{2016/08/13}{Bundle identifiers: 1) Added a line to the Makefile to pick out the date of the most-recently-updated file and put it in a file called VERSION (suggestion from Petra Rübe-Pugliese at CTAN based on their docs).; 2) Regularised the identity of the version date from a global variable date, rather than working it out each time..}
% \changes{v1.025}{2016/08/12}{README: Title and identity moved to a template and deleted from readme.xml.}
% \changes{v1.02}{2016/07/05}{Moved specification for babel to pre-options: 1) The babel package is now pre-specified with the PassOptionsToPackage command, to avoid conflicts with options later; 2) Now using \XeLaTeX{} and biber.}
% \changes{v1.01}{2016/06/12}{Edited out duplications in documentation: 1) Tidied explanations and documented the markup in more detail; 2) Now using \XeLaTeX{} and biber.}
% \changes{v1.0}{2016/06/12}{Tidied up the processing and greatly extended documentation: 1) Rewrote almost all the description of how to create the basic XML file; 2) The RFC2119 Note now automatically includes the relevant \BibTeX{} entry in the .bib file (done in the Makefile).}
% \changes{v0.92}{2016/05/12}{Major changes to the way the .dtx code is created: 1) Moved the specification of the processor from an attribute on the root element to the part start-tag for documentation; 2) Switched from the build file to a Makefile (still incomplete).; 3) Enabled the use of \XeLaTeX{} for the documentation in the .dtx file; 4) Largely rewrote the code for auto-detecting and adding packages, combining the auto-detect packages with those specified by the author. It now includes: a] negatable crosslinks specified in the remap on the constructorsynopsis start-tag (eg specifying the use of \XeLaTeX{} automatically invokes fontspec and defeats the autoload of the inputenc and fontenc packages), b] addition of the use of the arch and condition attributes on the bibliography start-tag to hold the \BibTeX{} package (biblatex) and the \BibTeX{} engine (biber), c] changes from the text-mode README to a MarkdownREADME.md to accommodate TUG’s submission process; 5) The layout of the zip files has been significantly changed to accord with the submission requirements of CTAN; 6) A Bibliography of the documentation is now extracted to disk directly in the XSLT3 process, rather than via the VerbatimOut environment in the .dtx file; 7) The RFC2119 Note now automatically includes the relevant \BibTeX{} entry in the .bib file (done in the Makefile).}
% \changes{v0.915}{2016/05/11}{Merged: Merged default packages with author-requested ones..}
% \changes{v0.91}{2016/05/10}{Regenerated: Recreated package with new classpack code to create zip file to the CTAN standard.}
% \changes{v0.905}{2015/10/06}{Enhanced: Added support for \XeLaTeX{}.}
% \changes{v0.9}{2015/12/19}{Removed fixltx2e: The code from the fixltx2e package has now been merged into the kernel, so the preload of this package is no longer needed..}
% \changes{v0.82}{2014/03/31}{Fixed hline bug: Should have been renewcommand not def.}
% \changes{v0.81}{2014/03/14}{Added more logo commands: XeTeX, XeLaTeX.}
% \changes{v0.80}{2014/03/10}{Moved doc commands from XSLT2: ToC settings and revmarg.}
% \changes{v0.79}{2014/03/07}{Editorial update: Small corrections.}
% \changes{v0.78}{2014/03/06}{Changed default code color: MacroFont now DarkBlue.}
% \changes{v0.77}{2013/05/28}{Removed unwanted definitions: classorpackage.}
% \changes{v0.76}{2013/05/05}{Modified documentation: Started working on Makefile.}
% \changes{v0.75}{2013/03/14}{Added secondary files: Secondary output files possible; reversed usage of role attribute on keywords;.}
% \changes{v0.74}{2013/02/21}{Added experimental autopackage: This implements automated package inclusion based on the markup used by the author..}
% \changes{v0.73}{2012/06/19}{Added readme.xml and db2plaintext.xsl: This implements dynamic README generation..}
% \changes{v0.72}{2012/02/11}{Wrote internal documentation: Created the classpack.xml template as an example..}
% \changes{v0.71}{2012/02/08}{First time this was used to document itself: The title element and subtitle element are now subsumed beneath the generated title in the output..}
%
% \GetFileInfo{classpack.dtx}
%
% \DoNotIndex{\@,\@@par,\@beginparpenalty,\@empty}
% \DoNotIndex{\@flushglue,\@gobble,\@input,\@makefnmark}
% \DoNotIndex{\@makeother,\@maketitle,\@namedef,\@ne}
% \DoNotIndex{\@spaces,\@tempa,\@tempb,\@tempswafalse}
% \DoNotIndex{\@tempswatrue,\@thanks,\@thefnmark,\@topnum}
% \DoNotIndex{\@@,\@elt,\@forloop,\@fortmp,\@gtempa}
% \DoNotIndex{\@totalleftmargin,\",\/,\@ifundefined,\@nil}
% \DoNotIndex{\@verbatim,\@vobeyspaces,\|,\~,\ ,\active}
% \DoNotIndex{\advance,\aftergroup,\begingroup,\bgroup}
% \DoNotIndex{\mathcal,\csname,\def,\documentstyle}
% \DoNotIndex{\dospecials,\edef,\egroup,\else,\endcsname}
% \DoNotIndex{\endgroup,\endinput,\endtrivlist}
% \DoNotIndex{\expandafter,\fi,\fnsymbol,\futurelet,\gdef}
% \DoNotIndex{\global,\hbox,\hss,\if,\if@inlabel}
% \DoNotIndex{\if@tempswa,\if@twocolumn,\ifcase,\ifcat}
% \DoNotIndex{\iffalse,\ifx,\ignorespaces,\index,\input}
% \DoNotIndex{\item,\jobname,\kern,\leavevmode,\leftskip}
% \DoNotIndex{\let,\llap,\lower,\m@ne,\next,\newpage}
% \DoNotIndex{\nobreak,\noexpand,\nonfrenchspacing}
% \DoNotIndex{\obeylines,\or,\protect,\raggedleft}
% \DoNotIndex{\rightskip,\rm,\sc,\setbox,\setcounter}
% \DoNotIndex{\small,\space,\string,\strut,\strutbox}
% \DoNotIndex{\thefootnote,\thispagestyle,\topmargin}
% \DoNotIndex{\trivlist,\tt,\twocolumn,\typeout,\vss,\vtop}
% \DoNotIndex{\xdef,\z@,\,,\@bsphack,\@esphack,\@noligs}
% \DoNotIndex{\@vobeyspaces,\@xverbatim,\`,\catcode,\end}
% \DoNotIndex{\escapechar,\frenchspacing,\glossary}
% \DoNotIndex{\hangindent,\hfil,\hfill,\hskip,\hspace,\ht}
% \DoNotIndex{\it,\langle,\leaders,\long,\makelabel}
% \DoNotIndex{\marginpar,\markboth,\mathcode,\mathsurround}
% \DoNotIndex{\mbox,\newcount,\newdimen,\newskip}
% \DoNotIndex{\nopagebreak,\parfillskip,\parindent}
% \DoNotIndex{\parskip,\penalty,\raise,\rangle,\section}
% \DoNotIndex{\setlength,\TeX,\topsep,\underline,\unskip}
% \DoNotIndex{\verb,\vskip,\vspace,\widetilde,\\,\%,\@date}
% \DoNotIndex{\@defpar,\[,\{,\},\],\count@,\ifnum,\loop}
% \DoNotIndex{\today,\uppercase,\uccode,\baselineskip}
% \DoNotIndex{\begin,\tw@,\a,\b,\c,\d,\e,\f,\g,\h,\i,\j,\k}
% \DoNotIndex{\l,\m,\n,\o,\p,\q,\r,\s,\t,\u,\v,\w,\x,\y,\z}
% \DoNotIndex{\A,\B,\C,\D,\E,\F,\G,\H,\I,\J,\K,\L,\M,\N,\O}
% \DoNotIndex{\P,\Q,\R,\S,\T,\U,\V,\W,\X,\Y,\Z,\1,\2,\3,\4}
% \DoNotIndex{\5,\6,\7,\8,\9,\0,\!,\#,\$,\&,\',\(,\)}
% \DoNotIndex{\+,\.,\:,\;,\<,\=,\>,\?,\_,\discretionary}
% \DoNotIndex{\immediate,\makeatletter,\makeatother}
% \DoNotIndex{\meaning,\newenvironment,\par,\relax}
% \DoNotIndex{\renewenvironment,\repeat,\scriptsize}
% \DoNotIndex{\selectfont,\the,\undefined,\arabic,\do}
% \DoNotIndex{\makeindex,\null,\number,\show,\write,\@ehc}
% \DoNotIndex{\@author,\@ehc,\@ifstar,\@sanitize,\@title}
% \DoNotIndex{\everypar,\if@minipage,\if@restonecol,\ifeof}
% \DoNotIndex{\ifmmode,\lccode,\newtoks,\onecolumn,\openin}
% \DoNotIndex{\p@,\SelfDocumenting,\settowidth}
% \DoNotIndex{\@resetonecoltrue,\@resetonecolfalse,\bf}
% \DoNotIndex{\clearpage,\closein,\lowercase}
% \DoNotIndex{\@inlabelfalse,\selectfont,\mathcode}
% \DoNotIndex{\newmathalphabet,\rmdefault,\bfdefault}
% \DoNotIndex{\DeclareRobustCommand}
% \DoNotIndex{\PassOptionsToPackage}
% \DoNotIndex{\MacroFont}
% \DoNotIndex{\subsubsection}
% \setcounter{secnumdepth}{5}
% \setcounter{tocdepth}{5}
% \setcounter{IndexColumns}{2}
% \setlength{\columnsep}{3pc}
%
% \def\fileversion{1.19}
% \def\filedate{2020/05/19}
% \title{The  \textsf{classpack} \LaTeXe\ package\thanks{%
% This document corresponds to \textsf{classpack}
% \textit{v.}\ \fileversion , dated \filedate.}
% \\[1em]\Large 
% XML mastering for \LaTeX{} classes and packages
% \\[1ex]\large 
% Literate-programming solution \\for class and
%       package maintenance}
% \author{Peter Flynn\\\normalsize Silmaril Consultants\\[-.25ex]\normalsize Textual Therapy Division\\\normalsize(\url{peter@silmaril.ie})}
% \maketitle
% \renewcommand{\abstractname}{Summary}\thispagestyle{empty}
% \begin{abstract}
% \parskip=0.5\baselineskip
% \advance\parskip by 0pt plus 2pt
% \parindent=0pt% \noindent
% \LaTeX{} document classes and packages are conventionally
% created, maintained, and distributed in Doc\TeX{}
% ({\ttfamily{}.dtx}) format using the 
% \textsf{ltxdoc} class, which provides for
% interleaved code and documentation (`literate
%   programming'). However, the accurate construction of
% these files is technically challenging, and editing them is
% tedious and error-prone.\par
% \emph{ClassPack} allows a developer to
% create a \emph{DocBook5} {\smaller XML} document for a class or package,
% containing the documentation and annotated code, and it
% provides {\smaller XSLT3} scripts to generate the
% {\ttfamily{}.dtx}, {\ttfamily{}.ins}, and
% other files, which can be combined into a zip file suitable
% for submission to {\smaller CTAN}.\par
% \emph{This package \textsf{classpack}
%   contains the small typographic adjustments and utilities
%   needed to re-typeset the documentation of classes and
%   packages developed using \textsf{classpack-dev}.
%   It is not used for any other purpose and is not required for
%   using any other class or package, only for the documentation
%   of \emph{ClassPack}–developed classes
%   and packages.}\par
% \textbf{You do not need to install the
%   upcoming development package \textsf{classpack-dev}
%   unless you want to develop classes or packages by using
%   \emph{ClassPack}
%   yourself.}\par
% \emph{ClassPack} is a
% work-in-progress. A paper describing an earlier version was
% presented at the Balisage markup conference in Montréal
% \parencite{flynn2013}.\par
% \end{abstract}
% \clearpage
% \tableofcontents
% \clearpage
% \section*{Latest changes}
% \subsection*{v.1.19 (2020-05-19)}
% \begin{itemize}
% \item Prefixed local variables with CPK\par
% \item CTAN candidate\par
% \end{itemize}
% \subsection*{v.1.18 (2020-04-01)}
% \begin{itemize}
% \item Split project into \textsf{classpack}
% and \textsf{classpack-dev}\par
% \item Replaced \texttt{@conformance}
% with \texttt{@YYYY-MM-DD} on all
% \texttt{date} elements in the revision history. This
% meant a complete regression run after editing all
% source masters to use this change\par
% \item Fixed bug in makechapapp where it failed to add
% the filetype before substringing the full name of the
% document\par
% \end{itemize}
% \subsection*{v.1.17 (2020-03-19)}
% \begin{itemize}
% \item Documentation nearly finished\par
% \item Added remaining XSLT files and scripts\par
% \item Recoded the distinction between appendixes in the
%       code part and chapters in the files part.\par
% \end{itemize}
% \subsection*{v.1.16 (2019-10-20)}
% \begin{itemize}
% \item Documentation ongoing\par
% \item Added detection of options on url and xcolor\par
% \item Changed order of packages in
% {\ttfamily{}prepost.xml} to prevent option
% clashes\par
% \end{itemize}
% See p.\thinspace\pageref{changehistory} for earlier changes.
% \clearpage
% \section{The \emph{ClassPack} package}\label{pkgintro}
% \emph{ClassPack} itself is an
% {\smaller XML}-based class and package development and
% management system described in the documentation for the
% \textsf{classpack-dev} package. This package,
% \textsf{classpack}, is responsible only for
% formatting the documentation of classes and packages developed
% using \emph{ClassPack}.\par
% The only operating \LaTeX{} code is therefore the settings
% and ancillary commands summarised in \vref{cpsum}
% and documented in \vref{pkgmacros}.\par
% \subsection{Features}\label{cpsum}
% Details of the annotated code are in \vref{pkgmacros}. They cover the following formatting
%   changes:\par
% \begin{enumerate}
% \item Two-column index instead of three-column\par
% \item Dark Blue colour for annotated code\par
% \item Recalculation of the left-hand margin in
%       documentation to accommodate long variable names\par
% \item Wider space for the section numbers and page numbers
%       in the Table of Contents, and ragged-right setting for
%       section titles to prevent hyphenation\par
% \item Appendixes (used for annotated code for ancillary
%       files) are styled at section level, not chapter
%       level\par
% \item Definitions for the Con\TeX t{}, \XeTeX{}, and
%       \XeLaTeX{} logos (borrowed from the TUGboat style)\par
% \item Fake small caps (also from TUGboat) for the \BibTeX{}
%       logo\par
% \item New struts for adjusting table spacing, and an arrow
%       between examples of menu items in documentation\par
% \item Some hyphenation oddities\par
% \item A fix for the broken description environment item
%       label to stop it overflowing\par
% \item A counter to be used when calculating the number of
%       items in a list\par
% \end{enumerate}
% \clearpage
% \raggedright
% \RaggedRight\markboth{\refname}{\refname}\printbibliography[heading=shortbib]
% \StopEventually{\label{endcode}%
%   \clearpage
%   \newgeometry{left=3cm}%
%   \addcontentsline{toc}{section}{Change History}%
%   \label{changehistory}%
%   \PrintChanges
%   \clearpage
%   \label{codeindex}%
%   \addcontentsline{toc}{section}{Index}%
%   \PrintIndex}
% \addtolength{\CPKrevmarg}{\widthof{\LabelFont{descriptionlabel}}}
% \newgeometry{left=\CPKrevmarg}
% \message{Margin reset to \the\CPKrevmarg, to fit <descriptionlabel>}
% \iffalse
%<*package>
% \fi
% \clearpage
% \section{The \textsf{classpack} macros and
% settings}\label{pkgmacros}
% \subsection{Auto-initialisation}\label{pkgmacros:autoinit}
% This section is added automatically by \textit{ClassPack} 
% as a preamble to all classes and style packages. 
% The \textsf{fixltx2e} package, which used to be included
% automatically, is no longer preloaded, as its
% features are now a part of the latest \LaTeXe\ kernel.\par
% The code starts with identity and requirements which are 
% generated automatically as needed by the Doc\TeX\ system.
% For details see the \textsf{ltxdoc} package documentation.
% \par\smallskip
% \begingroup\color{DarkRed}\tabcolsep3pt\footnotesize
% \begin{tabular}{>{\refstepcounter{CodelineNo}\tiny\theCodelineNo}r@{\enspace}l}
% &\verb`\NeedsTeXFormat{LaTeX2e}[2015/01/01]`\\
% &\verb`\ProvidesPackage{classpack}[2020/05/19 v1.19`\\
% &\verb`  Macros for ClassPack documentation]`
% \end{tabular}\endgroup
% \setcounter{CodelineNo}{3}
%\iffalse
%%
%% Packages required for the class or package
%% 
% \fi
% \subsection{Packages required for the package}\label{stypackages}
% \begin{CPK@package}{graphicx}
% Provide for graphics (PNG, JPG, or PDF format (only) for
% pdflatex; EPS format (only) for standard \LaTeX{}).
% \iffalse
%% Provide for graphics (PNG, JPG, or PDF format (only) for pdflatex; EPS format (only) for standard \LaTeX{}).
% \fi
%    \begin{macrocode}
\RequirePackage{graphicx}%
%    \end{macrocode}
% \end{CPK@package}
% \begin{CPK@package}{array}
% Additional column formatting types for tables.
% \iffalse
%% Additional column formatting types for tables.
% \fi
%    \begin{macrocode}
\RequirePackage{array}%
%    \end{macrocode}
% \end{CPK@package}
% \begin{CPK@package}{url}
% Handling of URI formatting.
% \iffalse
%% Handling of URI formatting.
% \fi
%    \begin{macrocode}
\RequirePackage{url}%
%    \end{macrocode}
% \end{CPK@package}
% \begin{CPK@package}{marginnote}
% Adds more flexibiltiy to marginal notes.
% \iffalse
%% Adds more flexibiltiy to marginal notes.
% \fi
%    \begin{macrocode}
\RequirePackage[fulladjust]{marginnote}%
%    \end{macrocode}
% \end{CPK@package}
% 
% \subsection{Index settings}\label{index}
% \iffalse
%% 
%% ****************************************************************** 
%% 
%% Index settings
% \fi
% \begin{CPK@counter}{IndexColumns}\label{counter--IndexColumns}
% The \textsf{doctex} package uses a default
%     three-column index for the documentation, which is too
%     narrow for most purposes. We therefore make the index in
%     two columns, and space them slightly farther apart. We
%     test first for the existence of the counter, in case this
%     gets used in a document other than a
%     {\ttfamily{}.dtx} file. No such test is needed for
%     \DescribeLength{columnsep}{\ttfamily{}columnsep} because it is
%     defined in the \LaTeX{} kernel.\par
%    \begin{macrocode}
\@ifundefined{c@IndexColumns}{}{\setcounter{IndexColumns}{2}}
\setlength{\columnsep}{3pc}
%    \end{macrocode}
% \end{CPK@counter}
% \subsection{Annotation settings}\label{annset}
% \iffalse
%% 
%% ****************************************************************** 
%% 
%% Annotation settings
% \fi
% \begin{CPK@macro}{\MacroFont}\label{macro--MacroFont}
% The \textsf{doc} and \textsf{docx}
%     packages use the {\ttfamily{}\textbackslash{}MacroFont} command for
%     the marginal labelling of code annotation. We redefine it
%     here to add the colour DarkBlue (from the
%     \textbf{\texttt{svgnames}} option to the
%     \textsf{xcolor} package).\par
%    \begin{macrocode}
\def\MacroFont{\fontencoding\encodingdefault
  \ttfamily\fontseries{m}\fontshape\updefault
  \small\selectfont\color{DarkBlue}}
%    \end{macrocode}
% \end{CPK@macro}
% \begin{CPK@length}{\CPKrevmarg}\label{length--CPKrevmarg}
% The default margin width is often not wide enough for
%     long macro names, so in the {\smaller XSLT3} code
%     in {\ttfamily{}db2dtx.xsl} we find the widest name
%     and add any excess over 25mm to the margin width. Here we
%     define and set the width parameter for this, which will get
%     reset later when calculated. The name has no at-sign, as
%     it operates in user mode.\par
%    \begin{macrocode}
\newlength{\CPKrevmarg}
\setlength{\CPKrevmarg}{25mm}
%    \end{macrocode}
% \end{CPK@length}
% \begin{CPK@macro}{\CPKrunningecho}\label{macro--CPKrunningecho}
% This allows alignment of the current annotation name
%     (from \texttt{@xreflabel}) as a
%     reminder in a marginal note in a fake subheading
%     implemented by a \texttt{bridgehead} element. Again, no
%     at-sign for a user-mode command.\par
%    \begin{macrocode}
\newcommand{\CPKrunningecho}[1]{\leavevmode
  \marginnote[\sloppy\raggedleft\color{LightGrey}\hspace{0pt}#1]%
             {\sloppy\raggedright\color{LightGrey}\hspace{0pt}#1}%
}
\let\marginfont\ttfamily
%    \end{macrocode}
% \end{CPK@macro}
% \subsection{Table of Contents}
% \iffalse
%% 
%% ****************************************************************** 
%% 
%% Table of Contents
% \fi
% \begin{CPK@macro}{\l@subsection}\label{macro--l@subsection}
% Documentation can sometimes have more than nine
%     subdivisions in sections, subsections, etc, and over 99
%     pages; and the default widths in the ToC are too narrow
%     for this, so we widen the space for the subsection number
%     by 0.4em:\par
%    \begin{macrocode}
\renewcommand*\l@subsection{%
  \@dottedtocline{2}{1.5em}{2.7em}}
%    \end{macrocode}
% \end{CPK@macro}
% \begin{CPK@macro}{\l@subsubsection}\label{macro--l@subsubsection}
% Similarly we increase the subsection number space by
%     0.4em, and its margin, so they align:\par
%    \begin{macrocode}
\renewcommand*\l@subsubsection{%
  \@dottedtocline{3}{4.2em}{3.6em}}
%    \end{macrocode}
% \end{CPK@macro}
% \begin{CPK@macro}{\@pnumwidth}\label{macro--@pnumwidth}
% The page number width is set to 3em instead of 1.55em:\par
%    \begin{macrocode}
\renewcommand{\@pnumwidth}{3em}
%    \end{macrocode}
% \end{CPK@macro}
% \begin{CPK@macro}{\@tocrmarg}\label{macro--@tocrmarg}
% And the right margin space goes up from 2.55em to 3em;
%   the addition of 1fil makes the section titles typeset
%   raggedright, so that hyphenation will not occur.\par
%    \begin{macrocode}
\renewcommand{\@tocrmarg}{4em plus1fil}
%    \end{macrocode}
% \end{CPK@macro}
% \subsection{Lower-level sectioning}
% \iffalse
%% 
%% ****************************************************************** 
%% 
%% Lower-level sectioning
% \fi
% \begin{CPK@macro}{\subsubsection}\label{macro--subsubsection}
% The {\ttfamily{}\textbackslash{}subsubsection} command is used
%     in bridgehead mode, so needs less space above and
%     below.\par
%    \begin{macrocode}
\renewcommand\subsubsection{%
  \@startsection{subsubsection}{3}{\z@}%
    {-1ex\@plus -.25ex \@minus -.25ex}%
    {1ex \@plus .25ex}%
    {\sffamily\normalsize\bfseries}}
%    \end{macrocode}
% \end{CPK@macro}
% \subsection{Appendix settings}
% \iffalse
%% 
%% ****************************************************************** 
%% 
%% Appendix settings
% \fi
% \begin{CPK@macro}{\appendix}\label{macro--appendix}
% Change the way the appendix command works so that
%     appendixes get section-type styling in documentation.\par
%    \begin{macrocode}
\renewcommand\appendix{\par
  \setcounter{section}{0}%
  \setcounter{subsection}{0}%
  \gdef\thesection{\@Alph\c@section}}
%    \end{macrocode}
% \end{CPK@macro}
% \subsection{\TeX{} and other logos}
% \iffalse
%% 
%% ****************************************************************** 
%% 
%% TeX and other logos
% \fi
% \TeX{} and \LaTeX{} are defined in the \LaTeX{} kernel, but
%   most of the others are not. The following definitions are
%   taken from the \textsf{ltugboat} package, used for
%   typesetting the TUGboat journal.\par
% \begin{CPK@macro}{\ConTeXt}\label{macro--ConTeXt}
% Con\TeX t{} is a typography and typesetting system meant
%     to provide users easy and consistent access to advanced
%     typographical control \parencite{wp-context}.\par
%    \begin{macrocode}
\def\ConTeXt{C\kern-.0333emon\-\kern-.0667em\TeX
  \kern-.0333emt}
%    \end{macrocode}
% \end{CPK@macro}
% \begin{CPK@macro}{\tubreflect}\label{macro--tubreflect}
% Borrow the reflection code from TUGboat.\par
%    \begin{macrocode}
\def\tubreflect#1{%
  \@ifundefined{reflectbox}{%
    \PackageError{classpack}%
      {A graphics package must be loaded for \string\XeTeX}%
      {Add the graphicx package to your Preamble}%
  }{% otherwise OK
    \ifdim \fontdimen1\font>0pt
      \raise 1.75ex \hbox{\kern.1em
        \rotatebox{180}{#1}}\kern-.1em
    \else
      \reflectbox{#1}%
    \fi
  }%
}
%    \end{macrocode}
% \end{CPK@macro}
% \begin{CPK@macro}{\tubhideheight}\label{macro--tubhideheight}
% Borrow the method of hiding the height from TUGboat
%     as well.\par
%    \begin{macrocode}
\def\tubhideheight#1{\setbox0=\hbox{#1}%
  \ht0=0pt \dp0=0pt \box0 }
%    \end{macrocode}
% \end{CPK@macro}
% \begin{CPK@macro}{\XeTeX}\label{macro--XeTeX}
% Define \XeTeX{} and \XeLaTeX{}.\par
%    \begin{macrocode}
\DeclareRobustCommand\Xe[1]{\leavevmode
  \tubhideheight{\hbox{X%
    \setbox0=\hbox{\TeX}\setbox1=\hbox{E}%
    \lower\dp0\hbox{\raise\dp1\hbox{%
      \kern-.125em\tubreflect{E}}}%
    \kern-.1667em #1}}}
\def\XeTeX{\Xe\TeX}
%    \end{macrocode}
% \end{CPK@macro}
% \begin{CPK@macro}{\XeLaTeX}\label{macro--XeLaTeX}
% Define \XeLaTeX{} using the existing macros.\par
%    \begin{macrocode}
\def\XeLaTeX{\Xe{\,\LaTeX}}
%    \end{macrocode}
% \end{CPK@macro}
% \begin{CPK@macro}{\SMC}\label{macro--SMC}
% Define a new small caps for use in \BibTeX{}, and an
%     error message to go with it (from the
%     \textsf{ltugboat} package).\par
%    \begin{macrocode}
\DeclareRobustCommand\SMC{%
  \ifx\@currsize\normalsize\small\else
   \ifx\@currsize\small\footnotesize\else
    \ifx\@currsize\footnotesize\scriptsize\else
     \ifx\@currsize\large\normalsize\else
      \ifx\@currsize\Large\large\else
       \ifx\@currsize\LARGE\Large\else
        \ifx\@currsize\scriptsize\tiny\else
         \ifx\@currsize\tiny\tiny\else
          \ifx\@currsize\huge\LARGE\else
           \ifx\@currsize\Huge\huge\else
            \small\SMC@unknown@warning
 \fi\fi\fi\fi\fi\fi\fi\fi\fi\fi
}
\newcommand\SMC@unknown@warning{\PackageError{classpack}%
  {\string\SMC: nonstandard text font size command -- using
  \string\small}%
  {Check the font size or scaling for \the\@currsize}}
\newcommand\textSMC[1]{{\SMC #1}}
%    \end{macrocode}
% \end{CPK@macro}
% \begin{CPK@macro}{\BIBTeX}\label{macro--BIBTeX}
% Finally, define \BibTeX{} in various forms.\par
%    \begin{macrocode}
	    \def\Bib{%
  \ifdim \fontdimen1\font>0pt
     B{\SMC\SMC IB}%
  \else
     \textsc{Bib}%
  \fi
}
\def\BibTeX{\Bib\kern-.08em \TeX}
\let\BiBTeX\BibTeX
\let\BIBTeX\BibTeX
%    \end{macrocode}
% \end{CPK@macro}
% \par\begingroup
%   \fboxsep1em\centering
%   \fbox{\begin{minipage}{0.8\columnwidth}\sffamily
%  \raggedright\parindent0pt
%  \parskip=.5\baselineskip
% \subsubsection*{\sffamily The \textsf{flexlogo} package}
% The \textsf{flexlogo} package (under
%   development) will make this section obsolete, as it allows
%   for the complete redefining of the \TeX{}, \LaTeX{}, and
%   related logos for non-CM fonts.\par
% \end{minipage}}\par\endgroup
% \subsection{Formatting additions}
% \iffalse
%% 
%% ****************************************************************** 
%% 
%% Formatting additions
% \fi
% \begin{CPK@macro}{\CPKvstrut}\label{macro--CPKvstrut}
% Define a strut that adjusts to the size of type, for
%   use in spacing table headers and footers.\par
%    \begin{macrocode}
\newcommand{\CPKvstrut}{\vrule height1.2em depth.6667ex width0pt}
%    \end{macrocode}
% \end{CPK@macro}
% \begin{CPK@macro}{\CPKmenusep}\label{macro--CPKmenusep}
% Define a macro to format an arrow between documentary
%     menu items. Probably no longer needed now that the
%     \textsf{menukeys} package is available\par
%    \begin{macrocode}
\def\CPKmenusep{\thinspace$\rightarrow$\thinspace\allowbreak}
%    \end{macrocode}
% \end{CPK@macro}
% \begin{CPK@macro}{\CPKprestrut}\label{macro--CPKprestrut}
% Also a strut to precede paragraph cells\dots{}\par
%    \begin{macrocode}
\newcommand{\CPKprestrut}{\vrule height1em width0pt}
%    \end{macrocode}
% \end{CPK@macro}
% \begin{CPK@macro}{\CPKpoststrut}\label{macro--CPKpoststrut}
% \dots{}and one to follow them.\par
%    \begin{macrocode}
\newcommand{\CPKpoststrut}{\vrule depth.5ex width0pt}
%    \end{macrocode}
% \end{CPK@macro}
% \begin{CPK@macro}{\hyphenation}\label{macro--hyphenation}
% Add some hyphenation oddities.\par
%    \begin{macrocode}
\hyphenation{ele-ment ele-ments attri-bute attri-butes 
  docu-ment docu-ments primi-tive helico-pter}
%    \end{macrocode}
% \end{CPK@macro}
% \begin{CPK@macro}{\descriptionlabel}\label{macro--descriptionlabel}
% Fix the broken description environment item
%     label. This also gets fixed in the
%     \textsf{enumitem} package so it will probably go
%     from here next version.\par
%    \begin{macrocode}
\renewcommand*\descriptionlabel[1]{%
  \hspace\labelsep\sffamily\bfseries #1}
%    \end{macrocode}
% \end{CPK@macro}
% \begin{CPK@counter}{CPKcoref}\label{counter--CPKcoref}
% Add the counter to enable the use of \texttt{coref}
%     list counters.\par
%    \begin{macrocode}
\newcounter{CPKcoref}
%    \end{macrocode}
% \end{CPK@counter}
% \iffalse
%</package>
% \fi
% \appendix
% \newgeometry{left=3cm}
% \clearpage
% \section{The \LaTeX{} Project Public License (v\thinspace{}1.3c)}\label{LPPL}
% \begin{quotation}\small\sffamily\parindent0pt
%   \parskip.5\baselineskip
%   \color{DarkBlue}\noindent
% Everyone is allowed to distribute verbatim copies of this
%       license document, but modification of it is not allowed.\par
% \end{quotation}
% \subsection{Preamble}\label{Preamble}
% The \LaTeX{} Project Public License ({\smaller LPPL})
%       is the primary license under which the \LaTeX{} kernel and the
%       base \LaTeX{} packages are distributed.\par
% You may use this license for any work of which you hold the
%       copyright and which you wish to distribute.  This license may be
%       particularly suitable if your work is \TeX{}-related (such as a
%       \LaTeX{} package), but it is written in such a way that you can
%       use it even if your work is unrelated to \TeX{}.\par
% The section “Whether and How to Distribute Works under This
%       License”, below, gives instructions, examples, and
%       recommendations for authors who are considering distributing
%       their works under this license.\par
% This license gives conditions under which a work may be
%       distributed and modified, as well as conditions under which
%       modified versions of that work may be distributed.\par
% We, the \LaTeX{3} Project, believe that the conditions below
%       give you the freedom to make and distribute modified versions of
%       your work that conform with whatever technical specifications
%       you wish while maintaining the availability, integrity, and
%       reliability of that work.  If you do not see how to achieve your
%       goal while meeting these conditions, then read the document
%       {\ttfamily{}cfgguide.tex} and {\ttfamily{}modguide.tex} in the base \LaTeX{}
%       distribution for suggestions.\par
% \subsection{Definitions}\label{Definitions}
% In this license document the following terms are used:\par
% \begin{description}[style=unboxed]
% \item[Work\thinspace:]Any work being distributed under this License.\par
% \item[Derived Work\thinspace:]Any work that under any applicable law is derived from
%     the Work.\par
% \item[Modification\thinspace:]Any procedure that produces a Derived Work under any
%     applicable law~--- for example, the production of a file
%     containing an original file associated with the Work or a
%     significant portion of such a file, either verbatim or
%     with modifications and/or translated into another
%     language.\par
% \item[Modify\thinspace:]To apply any procedure that produces a Derived Work
%     under any applicable law.\par
% \item[Distribution\thinspace:]Making copies of the Work available from one person to
%     another, in whole or in part.  Distribution includes (but
%     is not limited to) making any electronic components of the
%     Work accessible by file transfer protocols such as
%     {\smaller FTP} or {\smaller HTTP} or by
%     shared file systems such as Sun's Network File System
%     ({\smaller NFS}).\par
% \item[Compiled Work\thinspace:]A version of the Work that has been processed into a
%     form where it is directly usable on a computer system.
%     This processing may include using installation facilities
%     provided by the Work, transformations of the Work, copying
%     of components of the Work, or other activities.  Note that
%     modification of any installation facilities provided by
%     the Work constitutes modification of the Work.\par
% \item[Current Maintainer\thinspace:]A person or persons nominated as such within the Work.
%     If there is no such explicit nomination then it is the
%     `Copyright Holder' under any applicable
%     law.\par
% \item[Base Interpreter\thinspace:]A program or process that is normally needed for
%     running or interpreting a part or the whole of the
%     Work.\par
% A Base Interpreter may depend on external components
%     but these are not considered part of the Base Interpreter
%     provided that each external component clearly identifies
%     itself whenever it is used interactively.  Unless
%     explicitly specified when applying the license to the
%     Work, the only applicable Base Interpreter is a
%     `\LaTeX{}-Format' or in the case of files
%     belonging to the `\LaTeX{}-format' a program
%     implementing the `\TeX{} language'.\par
% \end{description}
% \subsection{Conditions on Distribution and Modification}\label{Conditions}
% \begin{enumerate}
% \item Activities other than distribution and/or modification
%   of the Work are not covered by this license; they are
%   outside its scope. In particular, the act of running the
%   Work is not restricted and no requirements are made
%   concerning any offers of support for the Work.\par
% \item \label{item-distribute}You may distribute a complete, unmodified copy of the
%   Work as you received it.  Distribution of only part of the
%   Work is considered modification of the Work, and no right to
%   distribute such a Derived Work may be assumed under the
%   terms of this clause.\par
% \item You may distribute a Compiled Work that has been
%   generated from a complete, unmodified copy of the Work as
%   distributed under Clause~item~\ref{item-distribute} above above, as
%   long as that Compiled Work is distributed in such a way that
%   the recipients may install the Compiled Work on their system
%   exactly as it would have been installed if they generated a
%   Compiled Work directly from the Work.\par
% \item \label{item-currmaint}If you are the Current Maintainer of the Work, you may,
%   without restriction, modify the Work, thus creating a
%   Derived Work.  You may also distribute the Derived Work
%   without restriction, including Compiled Works generated from
%   the Derived Work.  Derived Works distributed in this manner
%   by the Current Maintainer are considered to be updated
%   versions of the Work.\par
% \item If you are not the Current Maintainer of the Work, you
%   may modify your copy of the Work, thus creating a Derived
%   Work based on the Work, and compile this Derived Work, thus
%   creating a Compiled Work based on the Derived Work.\par
% \item \label{item-conditions}If you are not the Current Maintainer of the Work, you
%   may distribute a Derived Work provided the following
%   conditions are met for every component of the Work unless
%   that component clearly states in the copyright notice that
%   it is exempt from that condition.  Only the Current
%   Maintainer is allowed to add such statements of exemption to
%   a component of the Work.\par
% \begin{enumerate}
% \item If a component of this Derived Work can be a direct
%       replacement for a component of the Work when that
%       component is used with the Base Interpreter, then,
%       wherever this component of the Work identifies itself to
%       the user when used interactively with that Base
%       Interpreter, the replacement component of this Derived
%       Work clearly and unambiguously identifies itself as a
%       modified version of this component to the user when used
%       interactively with that Base Interpreter.\par
% \item Every component of the Derived Work contains
%       prominent notices detailing the nature of the changes to
%       that component, or a prominent reference to another file
%       that is distributed as part of the Derived Work and that
%       contains a complete and accurate log of the
%       changes.\par
% \item No information in the Derived Work implies that any
%       persons, including (but not limited to) the authors of
%       the original version of the Work, provide any support,
%       including (but not limited to) the reporting and
%       handling of errors, to recipients of the Derived Work
%       unless those persons have stated explicitly that they do
%       provide such support for the Derived Work.\par
% \item You distribute at least one of the following with
%       the Derived Work:\par
% \begin{enumerate}
% \item A complete, unmodified copy of the Work; if your
%   distribution of a modified component is made by
%   offering access to copy the modified component from
%   a designated place, then offering equivalent access
%   to copy the Work from the same or some similar place
%   meets this condition, even though third parties are
%   not compelled to copy the Work along with the
%   modified component;\par
% \item Information that is sufficient to obtain a
%   complete, unmodified copy of the Work.\par
% \end{enumerate}
% \end{enumerate}
% \item If you are not the Current Maintainer of the Work, you
%   may distribute a Compiled Work generated from a Derived
%   Work, as long as the Derived Work is distributed to all
%   recipients of the Compiled Work, and as long as the
%   conditions of Clause~item~\ref{item-conditions} above, above, are met
%   with regard to the Derived Work.\par
% \item The conditions above are not intended to prohibit, and
%   hence do not apply to, the modification, by any method, of
%   any component so that it becomes identical to an updated
%   version of that component of the Work as it is distributed
%   by the Current Maintainer under Clause~item~\ref{item-currmaint} above, above.\par
% \item Distribution of the Work or any Derived Work in an
%   alternative format, where the Work or that Derived Work (in
%   whole or in part) is then produced by applying some process
%   to that format, does not relax or nullify any sections of
%   this license as they pertain to the results of applying that
%   process.\par
% \item % \begin{enumerate}
% \item A Derived Work may be distributed under a different
%       license provided that license itself honors the
%       conditions listed in Clause~item~\ref{item-conditions} in the list on p.\thinspace\pageref{item-conditions} above, in
%       regard to the Work, though it does not have to honor the
%       rest of the conditions in this license.\par
% \item If a Derived Work is distributed under a different
%       license, that Derived Work must provide sufficient
%       documentation as part of itself to allow each recipient
%       of that Derived Work to honor the restrictions in
%       Clause~item~\ref{item-conditions} in the list on p.\thinspace\pageref{item-conditions} above, concerning
%       changes from the Work.\par
% \end{enumerate}
% \item This license places no restrictions on works that are
%   unrelated to the Work, nor does this license place any
%   restrictions on aggregating such works with the Work by any
%   means.\par
% \item Nothing in this license is intended to, or may be used
%   to, prevent complete compliance by all parties with all
%   applicable laws.\par
% \end{enumerate}
% \subsection{No Warranty}\label{Warranty}
% There is no warranty for the Work.  Except when otherwise
%       stated in writing, the Copyright Holder provides the Work
%       `as is', without warranty of any kind, either
%       expressed or implied, including, but not limited to, the implied
%       warranties of merchantability and fitness for a particular
%       purpose.  The entire risk as to the quality and performance of
%       the Work is with you.  Should the Work prove defective, you
%       assume the cost of all necessary servicing, repair, or
%       correction.\par
% In no event unless required by applicable law or agreed to
%       in writing will The Copyright Holder, or any author named in the
%       components of the Work, or any other party who may distribute
%       and/or modify the Work as permitted above, be liable to you for
%       damages, including any general, special, incidental or
%       consequential damages arising out of any use of the Work or out
%       of inability to use the Work (including, but not limited to,
%       loss of data, data being rendered inaccurate, or losses
%       sustained by anyone as a result of any failure of the Work to
%       operate with any other programs), even if the Copyright Holder
%       or said author or said other party has been advised of the
%       possibility of such damages.\par
% \subsection{Maintenance of The Work}\label{Maintenance}
% The Work has the status `author-maintained'
%       if the Copyright Holder explicitly and prominently states near
%       the primary copyright notice in the Work that the Work can only
%       be maintained by the Copyright Holder or simply that it is
%       `author-maintained'.\par
% The Work has the status `maintained' if there
%       is a Current Maintainer who has indicated in the Work that they
%       are willing to receive error reports for the Work (for example,
%       by supplying a valid e-mail address). It is not required for the
%       Current Maintainer to acknowledge or act upon these error
%       reports.\par
% The Work changes from status `maintained' to
%       `unmaintained' if there is no Current Maintainer,
%       or the person stated to be Current Maintainer of the work cannot
%       be reached through the indicated means of communication for a
%       period of six months, and there are no other significant signs
%       of active maintenance.\par
% You can become the Current Maintainer of the Work by
%       agreement with any existing Current Maintainer to take over this
%       role.\par
% If the Work is unmaintained, you can become the Current
%       Maintainer of the Work through the following steps:\par
% \begin{enumerate}
% \item Make a reasonable attempt to trace the Current
%   Maintainer (and the Copyright Holder, if the two differ)
%   through the means of an Internet or similar search.\par
% \item If this search is successful, then enquire whether the
%   Work is still maintained.\par
% \begin{enumerate}
% \item If it is being maintained, then ask the Current
%       Maintainer to update their communication data within one
%       month.\par
% \item \label{item-intention}If the search is unsuccessful or no action to resume
%       active maintenance is taken by the Current Maintainer,
%       then announce within the pertinent community your
%       intention to take over maintenance.  (If the Work is a
%       \LaTeX{} work, this could be done, for example, by
%       posting to \url{news:comp.text.tex}.)\par
% \end{enumerate}
% \item % \begin{enumerate}
% \item If the Current Maintainer is reachable and agrees to
%       pass maintenance of the Work to you, then this takes
%       effect immediately upon announcement.\par
% \item \label{item-announce}If the Current Maintainer is not reachable and the
%       Copyright Holder agrees that maintenance of the Work be
%       passed to you, then this takes effect immediately upon
%       announcement.\par
% \end{enumerate}
% \item \label{item-change}If you make an `intention announcement'
%   as described in~item~\ref{item-intention} in the list on p.\thinspace\pageref{item-intention} above and after three
%   months your intention is challenged neither by the Current
%   Maintainer nor by the Copyright Holder nor by other people,
%   then you may arrange for the Work to be changed so as to
%   name you as the (new) Current Maintainer.\par
% \item If the previously unreachable Current Maintainer becomes
%   reachable once more within three months of a change
%   completed under the terms of~item~\ref{item-announce} in the list on p.\thinspace\pageref{item-announce}
%   or~item~\ref{item-change} above, then that
%   Current
%   Maintainer must become or remain the Current Maintainer upon
%   request provided they then update their communication data
%   within one month.\par
% \end{enumerate}
% A change in the Current Maintainer does not, of itself,
%       alter the fact that the Work is distributed under the
%       {\smaller LPPL} license.\par
% If you become the Current Maintainer of the Work, you should
%       immediately provide, within the Work, a prominent and
%       unambiguous statement of your status as Current Maintainer.  You
%       should also announce your new status to the same pertinent
%       community as in~item~\ref{item-intention} in the list on p.\thinspace\pageref{item-intention}
%       above.\par
% \subsection{Whether and How to Distribute Works under This
%       License}\label{Distribute}
% This section contains important instructions, examples, and
%       recommendations for authors who are considering distributing
%       their works under this license.  These authors are addressed as
%       `you' in this section.\par
% \subsubsection{Choosing This License or Another License}\label{Choosing}
% If for any part of your work you want or need to use
% \emph{distribution} conditions that differ
% significantly from those in this license, then do not refer to
% this license anywhere in your work but, instead, distribute
% your work under a different license. You may use the text of
% this license as a model for your own license, but your license
% should not refer to the {\smaller LPPL} or otherwise
% give the impression that your work is distributed under the
% {\smaller LPPL}.\par
% The document {\ttfamily{}modguide.tex} in the base \LaTeX{}
% distribution explains the motivation behind the conditions of
% this license.  It explains, for example, why distributing
% \LaTeX{} under the {\smaller GNU} General Public
% License ({\smaller GPL}) was considered inappropriate.
% Even if your work is unrelated to \LaTeX{}, the discussion in
% {\ttfamily{}modguide.tex} may still be
% relevant, and authors intending to distribute their works
% under any license are encouraged to read it.\par
% \subsubsection{A Recommendation on Modification Without
% Distribution}\label{WithoutDistribution}
% It is wise never to modify a component of the Work, even
% for your own personal use, without also meeting the above
% conditions for distributing the modified component.  While you
% might intend that such modifications will never be
% distributed, often this will happen by accident~--- you may
% forget that you have modified that component; or it may not
% occur to you when allowing others to access the modified
% version that you are thus distributing it and violating the
% conditions of this license in ways that could have legal
% implications and, worse, cause problems for the community. It
% is therefore usually in your best interest to keep your copy
% of the Work identical with the public one.  Many works provide
% ways to control the behavior of that work without altering any
% of its licensed components.\par
% \subsubsection{How to Use This License}\label{HowTo}
% To use this license, place in each of the components of
% your work both an explicit copyright notice including your
% name and the year the work was authored and/or last
% substantially modified.  Include also a statement that the
% distribution and/or modification of that component is
% constrained by the conditions in this license.\par
% Here is an example of such a notice and statement:\par
% \iffalse
%<*ignore>
% \fi
\begin{lstlisting}[language={[LaTeX]TeX}]
%%% pig.dtx
%%% Copyright 2005 M. Y. Name
%%
%% This work may be distributed and/or modified under the
%% conditions of the LaTeX Project Public License, either version 1.3
%% of this license or (at your option) any later version.
%% The latest version of this license is in
%%   http://www.latex-project.org/lppl.txt
%% and version 1.3 or later is part of all distributions of LaTeX
%% version 2005/12/01 or later.
%%
%% This work has the LPPL maintenance status `maintained'.
%% 
%% The Current Maintainer of this work is M. Y. Name.
%%
%% This work consists of the files pig.dtx and pig.ins
%% and the derived file pig.sty.
\end{lstlisting}
% \iffalse
%</ignore>
% \fi
% Given such a notice and statement in a file, the
% conditions given in this license document would apply, with
% the `Work' referring to the three files
% {\ttfamily{}pig.dtx}, {\ttfamily{}pig.ins}, and {\ttfamily{}pig.sty} (the last being generated
% from {\ttfamily{}pig.dtx} using {\ttfamily{}pig.ins}), the `Base
%   Interpreter' referring to any
% `\LaTeX{}-Format', and both `Copyright
%   Holder' and `Current Maintainer'
% referring to the person
% M.~Y.~Name\index{!}.\par
% If you do not want the Maintenance section of
% {\smaller LPPL} to apply to your Work, change
% `maintained' above into
% `author-maintained'. However, we recommend that
% you use `maintained' as the Maintenance
% section was added in order to ensure that your Work remains
% useful to the community even when you can no longer maintain
% and support it yourself.\par
% \subsubsection{Derived Works That Are Not Replacements}\label{NotReplacements}
% Several clauses of the {\smaller LPPL} specify
% means to provide reliability and stability for the user
% community. They therefore concern themselves with the case
% that a Derived Work is intended to be used as a (compatible or
% incompatible) replacement of the original Work. If this is not
% the case (e.g., if a few lines of code are reused for a
% completely different task), then clauses 6b and 6d shall not
% apply.\par
% \subsubsection{Important Recommendations}\label{Recommendations}
% \paragraph[Defining What Constitutes the Work]{Defining What Constitutes the Work :}
% The {\smaller LPPL} requires that distributions
%   of the Work contain all the files of the Work.  It is
%   therefore important that you provide a way for the licensee
%   to determine which files constitute the Work.  This could,
%   for example, be achieved by explicitly listing all the files
%   of the Work near the copyright notice of each file or by
%   using a line such as:\par
% \iffalse
%<*ignore>
% \fi
\begin{lstlisting}[language={[LaTeX]TeX}]
%% This work consists of all files listed in manifest.txt.
\end{lstlisting}
% \iffalse
%</ignore>
% \fi
% in that place.  In the absence of an unequivocal list it
%   might be impossible for the licensee to determine what is
%   considered by you to comprise the Work and, in such a case,
%   the licensee would be entitled to make reasonable
%   conjectures as to which files comprise the Work.\par
% \Finale

